\section{Mechanism}
\label{sec:mechanism}

We state the mechanism as a hypothesis and give the behavioural evidence already in hand, then the causal experiments that are planned but not yet run.

\paragraph{Hypothesis (target rebinding).}
The \texttt{<Instruct>}/\texttt{<Query>}/\texttt{<Document>} convention has no parser; it is plain text interpolated by the developer.
A second well-formed (query, document) pair inside the content is, to the model, indistinguishable from the developer's, so the model evaluates the last pair it sees.
It is not persuaded to answer ``safe''; it answers a different question about a different document.

\paragraph{Behavioural evidence (from \cref{sec:attack}).}
Four observations constrain the mechanism.
(1)~A forged document alone clears 9\% but a forged question--document pair clears 100\%: a complete trained shape is needed, not more text.
(2)~The bypass rate tracks the trailing document's standalone harm (\cref{fig:groundtruth}): the model grades that document honestly.
(3)~Softening the question without supplying a benign decoy fails (0--9\%): the question's role is to open a new pair, not to soften the check.
(4)~Prepending fails and label-stripping fails: position and wording are not sufficient, labelled structure is.
Together these rule out instruction-following, recency and persuasion, and are consistent only with the model binding ``the document under review'' to the last labelled pair.
The same signature holds on gpt-oss-safeguard (\cref{sec:crossmodel}).

\paragraph{Causal experiments (planned).}
Behavioural evidence localizes the effect to the pair structure but does not locate where in the network the binding happens.
We will run, on the open-weight guards:
\begin{itemize}[leftmargin=1.5em,itemsep=2pt]
  \item \textbf{Decoy dose--response at scale.} Thousands of decoys spanning the full harm range, fitting the bypass-rate curve against the decoy's standalone score. Expected result: \tbd{fitted curve, $n$}.
  \item \textbf{Activation patching.} Swap the residual-stream representations of the real and forged pair and test whether the verdict follows the forged pair. Expected localization: \tbd{layer range and site}.
  \item \textbf{Linear probe.} Train a probe for ``which document is being graded'' across layers and report where it becomes decodable. Expected result: \tbd{probe accuracy by layer}.
  \item \textbf{Shortcut check.} Following the role-separation shortcut analysis of \citet{wang2025illusion}, test benign-first orderings to confirm the model binds to the last pair rather than a fixed position.
\end{itemize}
The probe is motivated by role-confusion findings that injected text occupies the representation of the role it imitates~\citep{ye2026prompt}; we test the analogue, that the forged \texttt{<Query>} is represented like the real one.
